% compile: xelatex AUN3_4095102.tex
\documentclass{mko3}
\begin{document}
\mkothreetitle

% ----- หมวดที่ 1 ข้อมูลทั่วไป (รูปแบบ มคอ.3 AUN-QA) -----
\coursecode{4095102}
\coursename{การเขียนโปรแกรมขั้นสูงสำหรับการเรียนรู้ของเครื่อง (Advanced Programming for Machine Learning)}
\program{หลักสูตรวิศวกรรมศาสตรมหาบัณฑิต สาขาวิชาวิศวกรรมคอมพิวเตอร์และปัญญาประดิษฐ์}
\coursecredits{3 (2-2-5)}
\coursegroup{วิชาเฉพาะด้านบังคับ}
\yearlevel{นักศึกษาปริญญาโท — ภาคต้น (ตามแผนการศึกษา)}
\semester{1 / 2568}
\instructor{ผศ.ดร.พิศิษฐ์ นาคใจ และ ผศ.ดร.พัชรี มณีรัตน์}
\contact{pisit.nakjai@uru.ac.th, patcharee.manee@uru.ac.th}
\coursedesc{การเขียนโปรแกรมด้วยภาษาโปรแกรมสมัยใหม่สำหรับการเรียนรู้ของเครื่อง การเตรียมและจัดการข้อมูล การทำความสะอาดข้อมูล การแปลงข้อมูล การจัดการข้อมูลที่ขาดหาย การแบ่งชุดข้อมูล การแสดงผลข้อมูลด้วยแผนภาพและการตีความ การใช้ไลบรารีสำหรับสร้างและปรับปรุงแบบจำลอง การทดสอบและประเมินประสิทธิภาพ การใช้ระบบควบคุมเวอร์ชัน และการนำแบบจำลองไปใช้งานในสภาพแวดล้อมจริง}
\courseinfotable
\begin{clodomaintable}
  \cloitem{CLO1}{อธิบาย หลักการของภาษาโปรแกรมสมัยใหม่และไลบรารีสำหรับการเรียนรู้ของเครื่อง ได้อย่างถูกต้อง}{Cognitive}{2}
  \cloitem{CLO2}{ประยุกต์ใช้ เทคนิคการจัดการข้อมูล ตั้งแต่การเตรียม ทำความสะอาด แปลง และแบ่งชุดข้อมูลสำหรับแบบจำลองการเรียนรู้ของเครื่อง}{Cognitive}{3}
  \cloitem{CLO3}{วิเคราะห์ ข้อมูลด้วยแผนภาพและการตีความ เพื่อเลือกใช้และปรับปรุงแบบจำลองที่เหมาะสม}{Cognitive}{4}
  \cloitem{CLO4}{แสดงออก ถึงความตระหนักในจริยธรรมและความรับผิดชอบต่อสังคม ในการพัฒนาโปรแกรมสำหรับการเรียนรู้ของเครื่อง}{Affective}{3}
  \cloitem{CLO5}{พัฒนา แบบจำลองการเรียนรู้ของเครื่อง โดยใช้ระบบควบคุมเวอร์ชันและนำไปใช้งานในสภาพแวดล้อมจริงได้อย่างมีประสิทธิภาพ}{Psychomotor}{3}
\end{clodomaintable}

\begin{cloplotable}{5}{ & \textbf{PLO1} & \textbf{PLO2} & \textbf{PLO3} & \textbf{PLO4} & \textbf{PLO5}}
  \textbf{CLO1} & X &   &   &   &   \\ \hline
  \textbf{CLO2} & X & X &   &   &   \\ \hline
  \textbf{CLO3} &   & X & X &   &   \\ \hline
  \textbf{CLO4} &   &   &   & X &   \\ \hline
  \textbf{CLO5} & X &   & X &   &   \\ \hline
\end{cloplotable}

\begin{bloomtable}
  \bloomrow{CLO1}{2}{}{}{}
  \bloomrow{CLO2}{3}{}{}{}
  \bloomrow{CLO3}{4}{}{}{}
  \bloomrow{CLO4}{}{3}{}{}
  \bloomrow{CLO5}{}{}{3}{3}
\end{bloomtable}

\begin{contentanalysistable}
  \contentrow{หลักการของภาษาโปรแกรมสมัยใหม่และไลบรารีสำหรับการเรียนรู้ของเครื่อง}{Cognitive 2}{อธิบาย}{ได้อย่างถูกต้อง (U) เพื่อนำไปประยุกต์ใช้ในการพัฒนาโปรแกรมสำหรับการเรียนรู้ของเครื่อง (CLO1) [PLO1]}{4}
  \contentrow{เทคนิคการจัดการข้อมูล}{Cognitive 3}{ประยุกต์ใช้}{ในการจัดการข้อมูลตั้งแต่การเตรียม ทำความสะอาด แปลง และแบ่งชุดข้อมูล (Ap) เพื่อใช้ในแบบจำลองการเรียนรู้ของเครื่อง (CLO2) [PLO1, PLO2]}{4}
  \contentrow{ข้อมูลด้วยแผนภาพและการตีความ}{Cognitive 4}{วิเคราะห์}{ข้อมูลผ่านการแสดงผลด้วยแผนภาพ (An) เพื่อตีความและเลือกใช้แบบจำลองที่เหมาะสม (CLO3) [PLO2, PLO3]}{4}
  \contentrow{ความตระหนักในจริยธรรมและความรับผิดชอบต่อสังคม}{Affective 3}{แสดงออก}{ถึงจริยธรรมและความรับผิดชอบ (Va) ในการพัฒนาโปรแกรมสำหรับการเรียนรู้ของเครื่อง (CLO4) [PLO4]}{4}
  \contentrow{แบบจำลองการเรียนรู้ของเครื่อง}{Psychomotor 3}{พัฒนา}{แบบจำลองโดยใช้ระบบควบคุมเวอร์ชัน (A) และนำไปใช้งานในสภาพแวดล้อมจริงได้อย่างมีประสิทธิภาพ (CLO5) [PLO1, PLO3]}{4}
\end{contentanalysistable}

\begin{courseactivities}
  \actitem{ศึกษาเอกสารและวิดีโอ ทดลองใช้ไลบรารีตามตัวอย่าง และอภิปรายหลักการ (CLO1)}
  \actitem{ฝึกปฏิบัติการจัดการข้อมูลจริง และทำงานกลุ่มวิเคราะห์ข้อมูล (CLO2)}
  \actitem{สร้างแผนภาพวิเคราะห์ข้อมูล ตีความผลการวิเคราะห์ และอภิปรายผล (CLO3)}
  \actitem{วิเคราะห์กรณีศึกษาด้านจริยธรรม อภิปรายกลุ่ม และเขียนสะท้อนความคิด (CLO4)}
  \actitem{พัฒนาโครงงาน ML ใช้ Git ในการทำงาน deploy model และนำเสนอผลงาน (CLO5)}
\end{courseactivities}

\begin{teachingtable}
  \teachrow{CLO1}{บรรยายทฤษฎีและหลักการ - สาธิตการใช้ไลบรารี - ยกตัวอย่างการประยุกต์ใช้ [IEL: ประสบการณ์เชิงบูรณาการ ลงมือปฏิบัติ พัฒนาท้องถิ่น]}{แบบทดสอบความเข้าใจ - การนำเสนอและอธิบายหลักการ - การตอบคำถาม}
  \teachrow{CLO2}{สาธิตเทคนิคการจัดการข้อมูล - workshop การเตรียมข้อมูล - วิเคราะห์กรณีศึกษา - ให้คำแนะนำระหว่างปฏิบัติ [IEL: ประสบการณ์เชิงบูรณาการ ลงมือปฏิบัติ พัฒนาท้องถิ่น]}{แบบฝึกหัดการจัดการข้อมูล - โปรเจคกลุ่ม - การนำเสนอผลการวิเคราะห์}
  \teachrow{CLO3}{สอนเทคนิคการวิเคราะห์ข้อมูล - สาธิตการสร้างแผนภาพ - อธิบายการตีความผล - ให้คำปรึกษาการวิเคราะห์ [IEL: ประสบการณ์เชิงบูรณาการ ลงมือปฏิบัติ พัฒนาท้องถิ่น]}{รายงานการวิเคราะห์ข้อมูล - คุณภาพของแผนภาพ - การนำเสนอและอภิปรายผล}
  \teachrow{CLO4}{อภิปรายประเด็นจริยธรรม - ยกตัวอย่างกรณีศึกษา - เชิญวิทยากรมาแลกเปลี่ยน [IEL: ประสบการณ์เชิงบูรณาการ ลงมือปฏิบัติ พัฒนาท้องถิ่น]}{รายงานวิเคราะห์จริยธรรม - การมีส่วนร่วมในการอภิปราย - พฤติกรรมการทำงาน}
  \teachrow{CLO5}{สาธิตการใช้ Git/GitHub - แนะนำการ deploy model - ให้คำปรึกษาโครงงาน [IEL: ประสบการณ์เชิงบูรณาการ ลงมือปฏิบัติ พัฒนาท้องถิ่น]}{คุณภาพโค้ดและ version control - ประสิทธิภาพของแบบจำลอง - การนำเสนอโครงงาน}
\end{teachingtable}

\begin{assessmenttable}
  \assessrow{สอบก่อนเรียน (ถ้ามี)}{---}{0}
  \assessrow{สอบระหว่างเรียน}{CLO1, CLO2, CLO3}{30}
  \assessrow{สอบปลายภาค}{CLO1, CLO2, CLO3, CLO4, CLO5}{40}
  \assessrow{รายงาน / โครงงาน / กิจกรรม}{CLO3, CLO4, CLO5}{30}
\end{assessmenttable}

\begin{weeklyplan}
  \weekrow{1}{Python และไลบรารี ML พื้นฐาน}{บรรยาย + Lab}{CLO1}{แบบทดสอบ}
  \weekrow{2}{การจัดการข้อมูล: อ่านและทำความสะอาด}{Workshop}{CLO2}{แบบฝึกหัด}
  \weekrow{3}{การแปลงและจัดการข้อมูลขาดหาย}{Lab}{CLO2}{รายงานย่อย}
  \weekrow{4}{การแบ่งชุด Train/Validation/Test}{Lab + อภิปราย}{CLO2}{แบบฝึกหัด}
  \weekrow{5}{การแสดงผลข้อมูล (EDA)}{Workshop}{CLO3}{สอบกลาง 1}
  \weekrow{6}{การตีความผลและเลือกแบบจำลอง}{บรรยาย + Lab}{CLO3}{รายงาน}
  \weekrow{7}{สร้างและฝึกแบบจำลองด้วย scikit-learn}{Lab}{CLO3, CLO5}{แบบฝึกหัด}
  \weekrow{8}{ประเมินประสิทธิภาพแบบจำลอง}{Lab + กรณีศึกษา}{CLO3}{สอบกลาง 2}
  \weekrow{9}{Git/GitHub สำหรับ ML}{Workshop}{CLO5}{แบบฝึกหัด}
  \weekrow{10}{จริยธรรมและความรับผิดชอบใน ML}{อภิปราย + วิทยากร}{CLO4}{บทความสะท้อนคิด}
  \weekrow{11}{ออกแบบโครงงานและกำหนด scope}{PBL}{CLO5}{ข้อเสนอโครงงาน}
  \weekrow{12--13}{พัฒนาโครงงานและ deploy}{Mentoring + Lab}{CLO5}{รายงานความก้าวหน้า}
  \weekrow{14}{ทดสอบและปรับปรุงระบบ}{Lab}{CLO5}{สาธิต}
  \weekrow{15}{นำเสนอโครงงานและสอบปลายภาค}{นำเสนอ + สอบ}{CLO1--CLO5}{สอบปลายภาค + โครงงาน}
\end{weeklyplan}

% ----- หมวดที่ 5 Rubric -----
\begin{rubricsection}
\begin{subrubric}{CLO1}{อธิบาย หลักการของภาษาโปรแกรมสมัยใหม่และไลบรารีสำหรับการเรียนรู้ของเครื่อง ได้อย่างถูกต้อง}
  \rubricrow{4 (ดีเยี่ยม)}{อธิบายหลักการได้อย่างถูกต้องลึกซึ้ง เชื่อมโยงทฤษฎีกับบริบทวิชาชีพและวรรณกรรมที่เกี่ยวข้องได้อย่างเหมาะสม}
  \rubricrow{3 (ดี)}{อธิบายหลักการได้ถูกต้องส่วนใหญ่ ครอบคลุมประเด็นสำคัญตามมาตรฐานหลักสูตรปริญญาโท}
  \rubricrow{2 (พอใช้)}{อธิบายได้บางส่วน ยังขาดความลึกหรือการเชื่อมโยงเชิงวิชาการ}
  \rubricrow{1 (ต้องปรับปรุง)}{ไม่สามารถอธิบายได้อย่างถูกต้องหรือไม่ครบถ้วน}
\end{subrubric}

\begin{subrubric}{CLO2}{ประยุกต์ใช้ เทคนิคการจัดการข้อมูล ตั้งแต่การเตรียม ทำความสะอาด แปลง และแบ่งชุดข้อมูลสำหรับแบบจำลองการเรียนรู้ของเครื่อง}
  \rubricrow{4 (ดีเยี่ยม)}{ประยุกต์ใช้ได้อย่างถูกต้องและเหมาะสม แก้ปัญหาซับซ้อนได้ด้วยเหตุผลและหลักฐานสนับสนุน}
  \rubricrow{3 (ดี)}{ประยุกต์ใช้ได้ถูกต้องส่วนใหญ่ ตรงตามวัตถุประสงค์ของงาน}
  \rubricrow{2 (พอใช้)}{ประยุกต์ใช้ได้บางส่วน ยังไม่เหมาะสมหรือไม่สมบูรณ์}
  \rubricrow{1 (ต้องปรับปรุง)}{ไม่สามารถประยุกต์ใช้ได้หรือใช้ผิดแนวทาง}
\end{subrubric}

\begin{subrubric}{CLO3}{วิเคราะห์ ข้อมูลด้วยแผนภาพและการตีความ เพื่อเลือกใช้และปรับปรุงแบบจำลองที่เหมาะสม}
  \rubricrow{4 (ดีเยี่ยม)}{วิเคราะห์เชิงวิพากษ์อย่างเป็นระบบ ใช้หลักฐาน/วรรณกรรมสนับสนุน สรุปและเสนอแนวทางได้ชัดเจน}
  \rubricrow{3 (ดี)}{วิเคราะห์ได้ถูกต้องส่วนใหญ่ มีเหตุผลประกอบเพียงพอตามมาตรฐานปริญญาโท}
  \rubricrow{2 (พอใช้)}{วิเคราะห์ได้ตื้น หรือยังขาดเหตุผล/หลักฐานสนับสนุน}
  \rubricrow{1 (ต้องปรับปรุง)}{วิเคราะห์ไม่ถูกต้องหรือไม่สามารถวิเคราะห์ได้}
\end{subrubric}

\begin{subrubric}{CLO4}{แสดงออก ถึงความตระหนักในจริยธรรมและความรับผิดชอบต่อสังคม ในการพัฒนาโปรแกรมสำหรับการเรียนรู้ของเครื่อง}
  \rubricrow{4 (ดีเยี่ยม)}{แสดงทัศนคติ/จริยธรรมวิชาชีพได้ชัดเจน สม่ำเสมอ และสะท้อนคิดอย่างลึกซึ้ง}
  \rubricrow{3 (ดี)}{แสดงออกได้เหมาะสมส่วนใหญ่ มีการสะท้อนคิด}
  \rubricrow{2 (พอใช้)}{แสดงออกได้บางครั้ง ยังไม่ชัดเจนหรือไม่สม่ำเสมอ}
  \rubricrow{1 (ต้องปรับปรุง)}{ไม่แสดงออกหรือแสดงออกไม่เหมาะสม}
\end{subrubric}

\begin{subrubric}{CLO5}{พัฒนา แบบจำลองการเรียนรู้ของเครื่อง โดยใช้ระบบควบคุมเวอร์ชันและนำไปใช้งานในสภาพแวดล้อมจริงได้อย่างมีประสิทธิภาพ}
  \rubricrow{4 (ดีเยี่ยม)}{ออกแบบ/พัฒนางานวิจัยหรือระบบได้สมบูรณ์ มีคุณภาพระดับบัณฑิตศึกษา นำไปใช้หรือเผยแพร่เชิงวิชาการได้}
  \rubricrow{3 (ดี)}{พัฒนาได้ดี ตรงวัตถุประสงค์และมาตรฐานที่ยอมรับได้}
  \rubricrow{2 (พอใช้)}{พัฒนาได้บางส่วน ยังไม่สมบูรณ์}
  \rubricrow{1 (ต้องปรับปรุง)}{ไม่สามารถพัฒนางานได้ตามที่กำหนด}
\end{subrubric}
\end{rubricsection}


\begin{learningmaterials}
  \matitem{เอกสารประกอบการบรรยายและคู่มือ Lab (Python, pandas, scikit-learn)}
  \matitem{ตำรา: Géron, \textit{Hands-On Machine Learning with Scikit-Learn, Keras, and TensorFlow}}
  \matitem{Jupyter Notebook, ชุดข้อมูลตัวอย่าง, Google Colab/Cloud}
  \matitem{เอกสาร Git/GitHub และแนวปฏิบัติ MLOps เบื้องต้น}
  \matitem{สื่อ LMS และบทความอ้างอิงจาก IEEE/ACM}
\end{learningmaterials}

\begin{supporttable}
  \supportrow{ห้องเรียน Smart Classroom}{1 ห้อง}{พร้อมใช้}{ทันสมัย}
  \supportrow{ห้องปฏิบัติการคอมพิวเตอร์}{1 ห้อง}{พร้อมใช้}{ทันสมัย}
  \supportrow{GPU Server / Cloud Computing Resources}{ครบถ้วน}{พร้อมใช้}{ทันสมัย}
  \supportrow{เอกสารประกอบการสอน, ตำรา, วารสารวิชาการ}{ครบถ้วน}{พร้อมใช้}{ทันสมัย}
\end{supporttable}

\signatureblock{ผศ.ดร.พิศิษฐ์ นาคใจ และ ผศ.ดร.พัชรี มณีรัตน์}{ผศ.ดร.พิทักษ์ คล้ายชม}

\end{document}